\documentclass[11pt,a4paper]{article}
\usepackage[italian]{babel}
\usepackage[margin=2.5cm]{geometry}
\usepackage{parskip}
\usepackage{amsmath, amssymb, amsthm}
\usepackage{booktabs}
\usepackage{array}
\usepackage{xcolor}
\usepackage{needspace}
\usepackage{enumitem}
\usepackage{listings}
\usepackage{tikz}
\usepackage[most]{tcolorbox}
\usepackage[hidelinks]{hyperref}
\usetikzlibrary{decorations.pathreplacing, shapes.geometric, arrows.meta, positioning}

% Niente vedove/orfane: i paragrafi non lasciano righe isolate a fine/inizio pagina
\widowpenalty=10000
\clubpenalty=10000

% Elenchi più compatti
\setlist{itemsep=2pt, parsep=0pt, topsep=4pt, partopsep=0pt}

% --- Riquadri con bordo nero, senza numero (si spezzano tra due pagine) ---
% Uso: \begin{definizione} ... \end{definizione}
%      \begin{teorema}[Nome] ... \end{teorema}  (il nome è facoltativo)
\tcbset{nota/.style={enhanced jigsaw, breakable, colback=white,
  colframe=black, boxrule=0.5pt, sharp corners}}
\NewTColorBox{definizione}{}{nota,
  before upper={\textbf{Definizione.}\ }}
\NewTColorBox{teorema}{o}{nota,
  before upper={\textbf{Teorema\IfValueT{#1}{ (#1)}.}\ }}
\NewTColorBox{proposizione}{o}{nota,
  before upper={\textbf{Proposizione\IfValueT{#1}{ (#1)}.}\ }}
\NewTColorBox{proprieta}{o}{nota,
  before upper={\textbf{Proprietà\IfValueT{#1}{ (#1)}.}\ }}
\NewTColorBox{assioma}{o}{nota, boxrule=1pt,
  before upper={\textbf{Assioma\IfValueT{#1}{ (#1)}.}\ }}

% --- Ambienti semplici (senza riquadro) ---
% Il titolo sta in cima al blocco, anche se il contenuto inizia con un riquadro o
% con codice; e non resta mai da solo in fondo a una pagina.
% Uso: \begin{esempio}[Titolo facoltativo] ... \end{esempio}
\newcommand{\newrunin}[2]{%
  \NewDocumentEnvironment{#1}{o}{%
    \par\Needspace{4\baselineskip}\addvspace{\medskipamount}\noindent
    \textbf{#2}\IfValueT{##1}{ (##1)}\textbf{.}\ \ignorespaces}%
  {\par\addvspace{\medskipamount}}}
\newrunin{esempio}{Esempio}
\newrunin{esercizio}{Esercizio}
\newrunin{osservazione}{Osservazione}

% V e F nelle tavole di verità
\newcommand{\V}{\textsf{V}}
\newcommand{\F}{\textsf{F}}

% --- Codice C ---
% Uso: \begin{codice} ... \end{codice}      codice C numerato
%      \begin{comando} ... \end{comando}    comandi da terminale
%      \begin{risultato} ... \end{risultato} output di un programma
%      \lstinline|printf("ciao");|           codice dentro al testo
\lstdefinestyle{c}{
  language=C,
  basicstyle=\ttfamily\small,
  keywordstyle=\color{blue}\bfseries,
  commentstyle=\color{gray}\itshape,
  stringstyle=\color{red!70!black},
  numbers=left,
  numberstyle=\tiny\color{gray},
  numbersep=8pt,
  columns=flexible,
  keepspaces=true,
  breaklines=true,
  showstringspaces=false,
  tabsize=4,
  morekeywords={include, define},
  % lettere accentate nei commenti
  literate={à}{{\`a}}1 {è}{{\`e}}1 {é}{{\'e}}1 {ì}{{\`i}}1
           {ò}{{\`o}}1 {ù}{{\`u}}1 {È}{{\`E}}1 {À}{{\`A}}1
}
\lstset{style=c}
\newtcblisting{codice}{listing only, nota, left=6mm,
  listing options={style=c}}
\newtcblisting{comando}{listing only, nota,
  listing options={basicstyle=\ttfamily\small, numbers=none, language={},
    columns=flexible, keepspaces=true, breaklines=true}}
\newtcblisting{risultato}{listing only, enhanced jigsaw, breakable,
  colback=black!5, colframe=black, boxrule=0.5pt, sharp corners,
  listing options={basicstyle=\ttfamily\small, numbers=none, language={},
    columns=flexible, keepspaces=true, breaklines=true}}

% --- Diagrammi di flusso (simboli standard) ---
\tikzset{
  terminale/.style = {ellipse, draw, minimum width=2.5cm, minimum height=0.9cm},
  io/.style        = {trapezium, trapezium left angle=70, trapezium right angle=110,
                      draw, minimum width=2.5cm, minimum height=0.9cm},
  processo/.style  = {rectangle, draw, minimum width=2.5cm, minimum height=0.9cm},
  decisione/.style = {diamond, draw, aspect=2, minimum width=2.5cm},
  freccia/.style   = {thick, -{Stealth}}
}

\title{Proprietà degli anelli e domini di integrità}
\author{Marco Cerbella}
\date{Lezione 10 -- 20 ottobre 2026}

\begin{document}
\maketitle

\section{Proprietà degli anelli}

Sia $(A, \oplus, \otimes)$ un anello e $a, b, c \in A$. Valgono le seguenti proprietà:
\begin{enumerate}
  \item $a \oplus b = a \oplus c \iff b = c$ (legge di cancellazione per $\oplus$);
  \item $a \otimes 0 = 0 \otimes a = 0$;
  \item $(-a)\otimes b = a \otimes (-b) = -(a \otimes b)$;
  \item $(-a)\otimes(-b) = a \otimes b$;
  \item $(a \ominus b)\otimes c = a \otimes c \ominus b \otimes c$ e $a \otimes (b \ominus c) = a \otimes b \ominus a \otimes c$ (dove $a \ominus b = a \oplus (-b)$);
  \item $(na)\otimes b = n\,(a \otimes b) = a \otimes (nb)$ per $n > 0$, dove $na = \underbrace{a \oplus a \oplus \dots \oplus a}_{n\text{ volte}}$;
  \item $(na)\otimes(mb) = (nm)(a \otimes b)$ per $n, m \in \mathbb{Z}$ (con $ka = \underbrace{(-a)\oplus\dots\oplus(-a)}_{|k|\text{ volte}}$ per $k$ negativo);
  \item se esiste l'unità $1 \in A$, allora $\forall a \in A$: $-a = (-1)\otimes a$.
\end{enumerate}

\section{Domini di integrità e divisori dello zero}

\begin{definizione}[Legge di annullamento del prodotto]
Sia $(A, \oplus, \otimes)$ un anello. Vale la \emph{legge di annullamento del prodotto} se
\[
  \forall a, b \in A \qquad a \otimes b = 0 \implies (a = 0 \ \text{oppure}\ b = 0).
\]
\end{definizione}

\begin{definizione}[Dominio di integrità]
Un anello $(A, \oplus, \otimes)$ è detto \emph{dominio di integrità} se vale la legge di annullamento del prodotto.
\end{definizione}

\begin{definizione}[Divisore dello zero]
Sia $(A, \oplus, \otimes)$ un anello. Un elemento $a \in A$ è detto \emph{divisore dello zero} se $a \neq 0$ e $\exists b \neq 0$, $b \in A$, tale che $a \otimes b = 0$ oppure $b \otimes a = 0$.
\end{definizione}

\begin{proposizione}
Sia $(A, \oplus, \otimes)$ un anello. Sono equivalenti:
\begin{enumerate}
  \item in $A$ vale la legge di annullamento del prodotto;
  \item in $A$ non ci sono divisori dello zero;
  \item l'insieme $A\setminus\{0\}$ è stabile per $\otimes$ (cioè $\otimes$ è un'operazione interna a $A\setminus\{0\}$).
\end{enumerate}
\end{proposizione}

\begin{proof}
$1) \Rightarrow 2)$. Per assurdo sia $a$ un divisore dello zero, cioè $a \neq 0$ ed esiste $b \neq 0$ tale che $a \otimes b = 0$. Per la legge di annullamento del prodotto allora $a = 0$ (impossibile) oppure $b = 0$ (impossibile): assurdo. Dunque non esistono divisori dello zero.

$2) \Rightarrow 3)$. Siano $a, b \in A\setminus\{0\}$; possiamo avere $a \otimes b = 0$? Se fosse $a \otimes b = 0$ allora $a \neq 0$, $b \neq 0$ e $a$ sarebbe un divisore dello zero, ma non esistono per ipotesi. Quindi $a \otimes b \neq 0$, cioè $a \otimes b \in A\setminus\{0\}$, cioè $\otimes$ è interna ad $A\setminus\{0\}$.

$3) \Rightarrow 1)$. Siano $a, b \in A$ tali che $a \otimes b = 0$. Per assurdo siano $a \neq 0$ e $b \neq 0$, cioè $a, b \in A\setminus\{0\}$; poiché $\otimes$ è interna a $A\setminus\{0\}$ si avrebbe $a \otimes b \neq 0$: impossibile. Quindi $a = 0$ oppure $b = 0$, cioè vale la legge di annullamento del prodotto.
\end{proof}

\begin{osservazione}
\begin{itemize}
  \item $(\mathbb{Z}, +, \cdot)$, $(\mathbb{Q}, +, \cdot)$, $(\mathbb{R}, +, \cdot)$ sono domini di integrità.
  \item In un campo (o in un corpo) vale la legge di annullamento del prodotto.
  \item In un dominio di integrità vale: se $a \neq 0$, allora $a \otimes b = a \otimes c \iff b = c$ (legge di cancellazione del prodotto).
\end{itemize}
\end{osservazione}

\begin{proof}[Dimostrazione della cancellazione]
$\Leftarrow$ Vale sempre. $\Rightarrow$ Da $a \otimes b = a \otimes c$ segue $0 = a \otimes c - a \otimes b = a \otimes (c - b)$. Per la legge di annullamento del prodotto $a = 0$ oppure $c - b = 0$; ma $a \neq 0$, quindi non può succedere la prima, dunque $c - b = 0$, cioè $c = b$.
\end{proof}

\subsection{Notazione per i campi}

Si indica un campo con $(\mathbb{K}, +, \cdot)$: $(\mathbb{K}, +, \cdot)$ è un anello commutativo con unità, e posto $\mathbb{K}^* = \mathbb{K}\setminus\{0\}$, ogni $k \in \mathbb{K}^*$ ha inverso $k^{-1}$. In un campo l'inverso moltiplicativo è unico.

\begin{esempio}[Divisori dello zero in $\mathbb{Z}_{12}$]
In $\mathbb{Z}_{12}$ la tavola del prodotto è
\[
\begin{array}{c|cccccc}
  \cdot & [1] & [2] & [3] & [4] & \cdots & [12] \\ \hline
  \relax[1] & [1] & [2] & [3] & [4] & \cdots & [12] \\
  \relax[2] & [2] & [4] & [6] & [8] & \cdots & [12] \\
  \relax[3] & [3] & [6] & [9] & [12] & \cdots & [12]
\end{array}
\]
$(\mathbb{Z}_{12}, +, \cdot)$ è un anello commutativo con unità; $[12]$ è l'elemento neutro per la somma, cioè $[0]$. Si ha $[3]\cdot[4] = [12] = [0]$, ma $[3] \neq [12]$ e $[4] \neq [12]$: non vale la legge di annullamento del prodotto e $[3]$, $[4]$ sono divisori dello zero.

Inoltre $[2]\cdot[1] = [2] = [2]\cdot[7]$ ma $[1] \neq [7]$: non vale la cancellazione del prodotto.
\end{esempio}

\end{document}
